% TOP_PROBLEM: {"id": 10100007, "title": "Limit shape of first-passage percolation", "category_id": 19, "category": {"id": 19, "name": "probability", "display_name": "Probability", "description": "Problems involving probability theory, stochastic processes, and random structures.", "slug": "probability", "order_index": 19, "created_at": "2026-07-31T15:26:25.671Z"}, "status": "open", "problem_number": "AMR-100-0007", "set_id": 13}
% FIRST_POSTED: 2026-10-01
% ENTRY_KIND: statement_only
% UnsolvedMath ID 10100007; source catalogue code AMR-100-0007
% !TeX program = lualatex
% Definition and sources only; this file is not an LLM solution attempt.
\documentclass[11pt,a4paper]{article}
\usepackage[margin=25mm]{geometry}
\usepackage{fontspec}
\setmainfont{lmroman10-regular.otf}[BoldFont=lmroman10-bold.otf,
 ItalicFont=lmroman10-italic.otf,BoldItalicFont=lmroman10-bolditalic.otf]
\usepackage{amsmath,amssymb,mathtools}
\usepackage{unicode-math}
\setmathfont{latinmodern-math.otf}
\AtBeginDocument{\let\setminus\smallsetminus}
\usepackage{xurl,hyperref}
\hypersetup{colorlinks=true,urlcolor=blue}
\setcounter{secnumdepth}{0}
\setlength{\parindent}{0pt}
\setlength{\parskip}{5pt}
\setlength{\emergencystretch}{3em}
\begin{document}
\section*{Limit shape of first-passage percolation}\label{10100007}
{\small UnsolvedMath ID 10100007 · AMR-100-0007 · Probability · AMR Open Problem Lists}
\subsection{Definitions and mathematical statement}
Fix $d\ge2$ and give $\mathbb Z^d$ its nearest-neighbor edges. Assign independent exponential passage times $\tau_e$ of mean one to all edges. For vertices $x,y$, define
\[
T(x,y)=\inf_{\gamma:x\to y}\sum_{e\in\gamma}\tau_e,
\]
where the infimum ranges over finite nearest-neighbor paths. Let $B(t)=\{x:T(0,x)\le t\}$, and replace each of its vertices by the unit cube centered there to obtain a subset $\widetilde B(t)\subset\mathbb R^d$.

The shape theorem supplies a deterministic compact convex body $K_d$ with nonempty interior such that $t^{-1}\widetilde B(t)$ converges to $K_d$ almost surely in Hausdorff distance. Equivalently, $K_d=\{x:\mu_d(x)\le1\}$, where the directional time constant is $\mu_d(x)=\lim_{n\to\infty}T(0,nx)/n$ for lattice directions and is extended to a norm on $\mathbb R^d$.

Determine this body explicitly, or equivalently determine the directional dependence of $\mu_d$. If only the shape up to scale is sought, normalize $K_d$ to Euclidean diameter one. Merely expressing the answer through the same unknown passage-time limit does not identify the shape.

The equivalent discrete growth process begins with only $0$ black, chooses uniformly among edges with exactly one black endpoint, and turns that edge's white endpoint black. Exponential memorylessness makes this the embedded jump chain of the passage-time model. Uniform choice is among boundary edges, so a white vertex with several black neighbors is proportionally more likely to be chosen. This distinction fixes the precise random growth model in the original question.
\subsection{Short English statement}
What is the deterministic macroscopic shape of exponential first-passage growth on a Euclidean lattice?
\subsection{Sources}
\begin{itemize}
\item[S1] ULAM.AI and the UnsolvedMath Contributors, supplied problems.json, record 10100007 (AMR-100-0007), AMR Open Problem Lists. Catalogue identity and source transcription; CC BY 4.0.
\url{https://huggingface.co/datasets/ulamai/UnsolvedMath}.
\item[S2] MathOverflow — Famous open problems in probability theory, answer 37167, source-order answer 7 of 13. Original source indexed by AMR-100-0007.
\url{https://mathoverflow.net/a/37167}.
\item[S3] B. Kloeckner, MathOverflow answer 37167 (2010), the original boundary-edge growth question.
\url{https://mathoverflow.net/a/37167}.
\item[S4] A. Auffinger, M. Damron and J. Hanson, 50 years of first passage percolation, shape theorem and limit-shape questions.
\url{https://arxiv.org/abs/1511.03262}.
\end{itemize}
\end{document}
